\documentclass[11pt]{article}
\usepackage{amsmath}
\usepackage{amssymb}
\usepackage{array}
\usepackage{geometry}
\geometry{margin=1in}
\begin{document}
\title{Unit conversion to a system with \(G=1\) for \sc{exp}}
\author{Martin D. Weinberg/draft}
\date{\today}
\maketitle
\section{Definitions}

Let the original (denoted as \emph{old}) numerical gravitational
constant be \(G_{\rm old}\). A particle in the old units has numeric
mass \(m\), position vector \(\mathbf{r}=(x,y,z)\) and velocity vector
\(\mathbf{v}=(u,v,w)\).  Define unit-scale factors as the number of
old units contained in one new unit (denoted as \emph{new}):
\begin{equation}
  s_L \equiv \frac{L_{\rm new}}{L_{\rm old}},\qquad
  s_V \equiv \frac{V_{\rm new}}{V_{\rm old}},\qquad
  s_M \equiv \frac{M_{\rm new}}{M_{\rm old}}.
  \label{eq:scalefactors}
\end{equation}
A numeric quantity in the new units (denoted with a prime) is obtained
by dividing the old numeric value by the corresponding scale factor:
\begin{equation}
  \mathbf{r}'=\frac{\mathbf{r}}{s_L},\qquad
  \mathbf{v}'=\frac{\mathbf{v}}{s_V},\qquad
  m'=\frac{m}{s_M}.
  \label{eq:scaledunits}
\end{equation}
Time-scale factor:
\[
  s_T \equiv \frac{T_{\rm new}}{T_{\rm old}}=\frac{s_L}{s_V}.
\]

\section{Dimensional constraint for \(G\) and the core relation}

Dimensional analysis gives \([G]=L\,V^2/M\). Thus numeric values
transform as
\[
  G_{\rm new} = G_{\rm old}\,\frac{s_M}{s_L\,s_V^2}.
\]
Requiring \(G_{\rm new}=1\) gives the central constraint
\[\boxed{ \; s_M = s_L\,s_V^{2}/G_{\rm old}\; }.\]

These formula give my usual recommendation to Gadget users: keep
length, velocity, and time units the same and scale your mass by
\(G_{old}\) (see eq. \ref{eq:scaledunits}).

\section{Conversion formulas (apply to each particle)}

Given \((m,\mathbf{r},\mathbf{v})\) in old units:
\[\boxed{ \; m'=\frac{m}{s_M},\quad \mathbf{r}'=\frac{\mathbf{r}}{s_L},\quad \mathbf{v}'=\frac{\mathbf{v}}{s_V},\quad t'=\frac{t}{s_T}=\frac{t\,s_V}{s_L}\; }.\]

\section{Solve-for formulas (pick any two scales, compute the third)}

From \(s_M = s_L s_V^2 / G_{\rm old}\):
\begin{itemize}
\item Given \(s_L\) and \(s_V\):
  \[s_M = \frac{s_L\,s_V^2}{G_{\rm old}}.\]
\item Given \(s_L\) and \(s_M\):
  \[s_V = \sqrt{\frac{s_M G_{\rm old}}{s_L}}.\]
\item Given \(s_V\) and \(s_M\):
\[s_L = \frac{s_M G_{\rm old}}{s_V^2}.\]
\end{itemize}
\section{Monopole (circular velocity) condition}

Suppose an enclosed mass (in old units) is \(M_{\rm enc}\) inside
radius \(r_0\) (old units). The old circular speed satisfies
\[
  v_{c,{\rm old}}^2=\frac{G_{\rm old}\,M_{\rm enc}}{r_0}.
\]
In the new units (where \(G_{\rm new}=1\)), at the same physical
radius (numerically \(r_0' = r_0/s_L\)) the circular speed squared
becomes
\[
  v_c'^2 = \frac{(M_{\rm enc}/s_M)}{(r_0/s_L)} = \frac{M_{\rm
      enc}\,s_L}{s_M\,r_0}.
\]
Using \(s_M=s_L s_V^2/G_{\rm old}\) this simplifies to
\[
  \boxed{ \; v_c'^2 = \frac{G_{\rm old}\,M_{\rm enc}}{s_V^2\,r_0} \; }
\]
or equivalently
\[
  \boxed{ \; s_V = \sqrt{\frac{G_{\rm old}\,M_{\rm enc}}{v_c'^2\,r_0}}\;. }
\]
Hence: if you demand a specific numeric circular speed value \(v_c'\)
at physical radius \(r_0\), this uniquely determines \(s_V\)
(independent of \(s_L\)). Once \(s_V\) is set, \(s_M\) follows from
\(s_M=s_L s_V^2/G_{\rm old}\); if you additionally set a target
numeric radius \(r_0'\) for \(r_0\), choose \(s_L = r_0/r_0'\).

\section{Summary table}

\begin{center}
\renewcommand{\arraystretch}{2}
\begin{tabular}{p{0.3\textwidth} p{0.65\textwidth}}
\hline
Task & Formula\\
\hline
  Core constraint & $s_M = \dfrac{s_L\,s_V^2}{G_{\rm old}}$ \\
  Convert variables & $r' = r/s_L,\quad v'=v/s_V,\quad m'=m/s_M,\quad t' = t\,s_V/s_L$ \\
  Solve for $s_M$ & $s_M = \dfrac{s_L\,s_V^2}{G_{\rm old}}$ \\
  Solve for $s_V$ & $s_V = \sqrt{\dfrac{G_{\rm old}\,s_M}{s_L}}$ \\
  Solve for $s_L$ & $s_L = \dfrac{G_{\rm old}\,s_M}{s_V^2}$ \\
  Circular-velocity condition & $s_V = \sqrt{\dfrac{G_{\rm old}\,M_{\rm enc}}{v_c'^2\,r_0}}$ \\
  Resulting $s_M$ & $s_M = \dfrac{s_L\,s_V^2}{G_{\rm old}} = s_L\,\dfrac{M_{\rm enc}}{v_c'^2\,r_0}$ \\
\hline
\end{tabular}
\end{center}

\section{Usage recipes}
\begin{itemize}
\item Pick new length and velocity units directly: specify \(s_L\) and
  \(s_V\), compute \(s_M = s_L s_V^2/G_{\rm old}\), then rescale
  particle numbers with the conversion formulas.
\item If you want \(v_c'=v_*\) at radius \(r_0\) for enclosed old mass
  \(M_{\rm enc}\): compute \(s_V\) from the circular-velocity formula,
  pick \(s_L\) (e.g. \(s_L=r_0\) if you want \(r_0'=1\)), then \(s_M\)
  follows.
\item If you require a particular old mass \(M_*\) to become \(m'=1\):
  set \(s_M=M_*\), then choose \(s_L\) (or \(s_V\)) and compute the
  remaining scale factor.
\end{itemize}
\end{document}